\begin{enumerate}
    \item On donne les couples d'oxydoréduction~: \;
          $\ce{I_{2}_{(aq)}}\ttslash{}\ce{I^{-}_{(aq)}}$ \; et \;
          $\ce{S4O^{2-}_{6(aq)}}\ttslash{}\ce{S2O^{2-}_{3(aq)}}.$

          Écrire la réaction d'oxydation des ions thiosulfate par le diiode en
          solution aqueuse.
    \item Cette réaction peut être utilisée pour doser le diiode en solution
          aqueuse.
        
          Dans un bécher, on place $\qty{20.0}{\mL}$ d'une solution aqueuse de
          diiode et on fait couler la solution d'ions thiosulfate contenue dans
          une burette graduée où la concentration molaire des ions thiosulfate
          est $\qty{0.10}{\mol\per\L}$.
    
          Dans le bécher, la solution jaune de diiode se décolore
          progressivement.
    
          Lorsqu'elle apparaît décolorée, on ajoute de l'empois d'amidon qui
          prend une couleur bleu foncé, révélant une petite quantité de diiode
          encore présente. On poursuit l'ajout de solution de thiosulfate
          jusqu'à décoloration de l'empois d'amidon. On a alors versé un volume
          $V=\qty{12.6}{\mL}$ de solution d'ions thiosulfate.

          On note $C$ la concentration molaire de la solution de diiode.
          \begin{enumerate}
              \item Dresser le tableau permettant de suivre l'évolution de la
                    transformation.
              \item Définir l'équivalence. Quel est l'avancement à l'équivalence
                    $x_{\eq}$~?
              \item Calculer la concentration $C$ de la solution de diiode.
          \end{enumerate}
\end{enumerate}


